\chapter{Reference: exit codes and cleanup}
\label{ch:refman-exit}

\begin{aside}
\code{run} returns an exit code. Use it. Distinguish
"user quit", "error", "interrupted", "success".
\end{aside}

\section{Working example}

\begin{lstlisting}[language=C++]
#include <maya/maya.hpp>

using namespace maya;
using namespace maya::dsl;

struct Exit {
    struct Model { int code = 0; };
    struct OkQuit {};
    struct ErrQuit {};
    using Msg = std::variant<OkQuit, ErrQuit>;

    static Model init() { return {}; }

    static auto update(Model m, Msg msg)
        -> std::pair<Model, Cmd<Msg>>
    {
        return std::visit(overload{
            [&](OkQuit) {
                m.code = 0;
                return std::pair{m, Cmd<Msg>::quit()};
            },
            [&](ErrQuit) {
                m.code = 2;
                return std::pair{m, Cmd<Msg>::quit()};
            },
        }, msg);
    }

    static Element view(const Model&) {
        return v(
            t<"choose an exit code"> | Bold,
            blank_,
            t<"q  exit 0 (ok)"> | Dim,
            t<"e  exit 2 (error)"> | Dim
        ) | pad<1> | border_<Round>;
    }

    static auto subscribe(const Model&) -> Sub<Msg> {
        return key_map<Msg>({
            {'q', OkQuit{}},
            {'e', ErrQuit{}},
        });
    }
};

int main() {
    Exit::Model final;
    // The runtime returns the exit code via run<>().
    // To propagate the Model-stored code, you can capture
    // it via a global, or return it from run() if your
    // runtime version supports model-on-quit.
    int rc = run<Exit>({.title = "exit"});
    return rc;
}
\end{lstlisting}

\section{Conventional codes}

\begin{itemize}[leftmargin=*]
  \item \code{0} — success.
  \item \code{1} — generic failure.
  \item \code{2} — bad usage.
  \item \code{130} — interrupted by Ctrl-C (\code{SIGINT}).
  \item \code{143} — killed by \code{SIGTERM}.
\end{itemize}

\section{Cleanup}

\begin{itemize}[leftmargin=*]
  \item RAII handles most of it; destructors run as
    \code{main} returns.
  \item For shared state, run cleanup inside a \code{Quit}
    handler before returning the \code{quit()} command.
  \item Don't restore the terminal yourself; \maya{} does it.
\end{itemize}

\section{Recap}

\begin{recap}
\begin{itemize}[leftmargin=*]
  \item \code{run} returns the exit code.
  \item Use conventional codes; scripts depend on them.
  \item RAII handles cleanup; trust it.
\end{itemize}
\end{recap}

\section*{Workshop}

\begin{enumerate}[leftmargin=*]
  \item Add three exit paths to your app: ok, cancel,
    error.
  \item Write a shell script that branches on the exit
    code.
  \item Verify cleanup runs after a forced \code{SIGTERM}.
\end{enumerate}
